% Level 2, Week 8 -- Testing: Prove It Works (homework)
\homeworkweek{8}{Testing: Prove It Works}{15--20 minutes. Do Core first; Challenge and Creative are optional.}
\namedateline

\begin{core}[Test the Cipher (about 10 minutes)]
Folder: \texttt{exercises/week08/hw-test-the-cipher}. \texttt{encode()} in \texttt{cipher\_tools.py} is the Caesar cipher from week 5, and it is correct.
Your job: write tests so nobody can ever break it without you noticing. Keep your paper cipher wheel nearby.
\begin{enumerate}[leftmargin=1.4em]
  \item Read \texttt{cipher\_tools.py}. Then open \texttt{starter.py} (your test file) and press \menu{Run}.
  \item TODO 1: work out \texttt{encode("xyz", 3)} with the wheel and write a test for it: \rule{3cm}{0.4pt}
  \item TODO 2: test a message with a space and a \texttt{!} in it. My message and its secret version:
  \answerlines{1}
  \item TODO 3: one more test of your own (ideas: shift 0, or an empty message \texttt{""}).
  \item Press \menu{Run}: all \texttt{PASSED}? Then run \texttt{check\_my\_work.py}. It tries your tests on three broken ciphers.
\end{enumerate}
\end{core}

\begin{challenge}[Spot the edge]
For each function, write one \textbf{normal} case and one \textbf{edge} case (no computer needed).
\begin{center}
\renewcommand{\arraystretch}{2.2}
\begin{tabular}{@{}>{\raggedright\arraybackslash}p{6.4cm}p{3.6cm}p{3.6cm}@{}}
\toprule
\textbf{Function} & \textbf{Normal case} & \textbf{Edge case} \\ \midrule
\texttt{count\_vowels(word)}: counts a, e, i, o, u & & \\ \midrule
\texttt{lives\_left(lives, hits)}: never below 0 & & \\ \midrule
\texttt{last\_item(items)}: the last item & & \\
\bottomrule
\end{tabular}
\end{center}
\end{challenge}

\begin{creative}[Test something of yours]
Pick a function you wrote in an earlier week (for example from your Monster Stats, Word Guess or Idea Notebook).
In the same folder, make a new file called \texttt{my\_tests.py}. Import your function (\py|from starter import my_function|), write two tests
(one normal, one edge), and copy the last three lines of this week's \texttt{starter.py} to the bottom so that \menu{Run} starts pytest.
This is exactly how your final project's tests will work.
Or: draw a comic strip of a \enquote{crash-test dummy} testing a game's strangest corner.
\end{creative}

\begin{selfcheck}
  \item My test names all start with \texttt{test\_} and each has an \texttt{assert}.
  \item All my tests say \texttt{PASSED} on the correct \texttt{encode()}.
  \item \texttt{check\_my\_work.py} says my tests catch all three broken ciphers.
  \item I can explain the difference between a normal case and an edge case.
\end{selfcheck}
